\documentclass[13pt,a4paper]{extarticle}
\usepackage{fontspec}
\usepackage{unicode-math}
\setmainfont{Latin Modern Roman}
\setmathfont{Latin Modern Math}
\usepackage{amsmath}
\usepackage[margin=1.8cm]{geometry}
\usepackage{enumitem}
\usepackage{array}
\usepackage{multicol}
\setlength{\columnsep}{1.2em}
\usepackage{eurosym}
\linespread{1.04}
\renewcommand{\arraystretch}{1.28}
\usepackage{newunicodechar}
\newunicodechar{·}{\kern-.13em\textperiodcentered\kern-.13em}
\usepackage{parskip}
\setlength{\parindent}{0pt}
\begin{document}
\section*{Exercicis i problemes}

\begin{quote}\small Llistat d'exercicis per practicar i consolidar els continguts del bloc. Segueix l'ordre dels apunts i inclou exercicis de diferents nivells de dificultat.\end{quote}



\vspace{1mm}\hrule\vspace{3mm}

\section*{1. Equacions de primer i segon grau}

\par\noindent\begin{minipage}{\linewidth}
\textbf{1.} Resol les equacions de primer grau següents:\par

\begin{tabular}{@{}p{.46\linewidth}p{.46\linewidth}@{}}
\textbf{a)} $3(x-2)-2(x+1)=5-x$ & \textbf{b)} $\dfrac{x-1}{2}-\dfrac{x+3}{3}=1-\dfrac{x}{6}$ \\[3pt]
\textbf{c)} $2(x+3)-x=x+6$ & \textbf{d)} $4x-2(2x-1)=3$ \\[3pt]
\end{tabular}\par
\end{minipage}\par\vspace{1.6mm}
\par\noindent\begin{minipage}{\linewidth}
\textbf{2.} Resol aquestes equacions de segon grau incompletes sense aplicar la fórmula:\par

\begin{tabular}{@{}p{.46\linewidth}p{.46\linewidth}@{}}
\textbf{a)} $3x^2-27=0$ & \textbf{b)} $5x^2+20x=0$ \\[3pt]
\textbf{c)} $2x^2+8=0$ & \textbf{d)} $x^2=7x$ \\[3pt]
\end{tabular}\par
\end{minipage}\par\vspace{1.6mm}
\par\noindent\begin{minipage}{\linewidth}
\textbf{3.} Resol les equacions de segon grau següents i indica, a partir del discriminant, quantes solucions té cadascuna:\par

\begin{tabular}{@{}p{.46\linewidth}p{.46\linewidth}@{}}
\textbf{a)} $x^2-5x+6=0$ & \textbf{b)} $4x^2-12x+9=0$ \\[3pt]
\textbf{c)} $x^2+x+1=0$ & \textbf{d)} $6x^2-x-2=0$ \\[3pt]
\end{tabular}\par
\end{minipage}\par\vspace{1.6mm}
\par\noindent\begin{minipage}{\linewidth}
\textbf{4.} Resol, simplificant primer les expressions:\par

\begin{itemize}[leftmargin=0pt,itemsep=6pt,topsep=4pt,parsep=0pt,label={}]
\item \textbf{a)} $(x-3)^2+(x+1)^2=2(x^2-7)$
\item \textbf{b)} $\dfrac{x(x+2)}{3}-\dfrac{x-1}{2}=\dfrac{x+5}{6}$
\item \textbf{c)} $(2x-1)(x+3)-(x+2)^2=3$
\item \textbf{d)} $0{,}5(x-1)^2-0{,}25(x+1)^2=4-x$
\end{itemize}
\end{minipage}\par\vspace{1.6mm}
\par\noindent\begin{minipage}{\linewidth}
\textbf{5.} Considera l'equació $x^2-6x+k=0$.\par

\begin{itemize}[leftmargin=0pt,itemsep=6pt,topsep=4pt,parsep=0pt,label={}]
\item \textbf{a)} Troba $k$ perquè tingui una solució doble i calcula-la.
\item \textbf{b)} Troba $k$ perquè $x=2$ sigui solució i calcula l'altra solució.
\item \textbf{c)} Per a quins valors de $k$ l'equació no té solució?
\end{itemize}
\end{minipage}\par\vspace{1.6mm}
\par\noindent\begin{minipage}{\linewidth}
\textbf{6.} Escriu una equació de segon grau amb coeficients enters que tingui com a solucions:\par

\begin{itemize}[leftmargin=0pt,itemsep=6pt,topsep=4pt,parsep=0pt,label={}]
\item \textbf{a)} $x=3$ i $x=-5$.
\item \textbf{b)} només $x=\dfrac23$ (solució doble).
\end{itemize}
\end{minipage}\par\vspace{1.6mm}
\par\noindent\begin{minipage}{\linewidth}
\textbf{7.} Resol les equacions següents, on $a$ i $b$ són constants:\par

\begin{itemize}[leftmargin=0pt,itemsep=6pt,topsep=4pt,parsep=0pt,label={}]
\item \textbf{a)} $ax+b=bx+a$, amb $a\neq b$
\item \textbf{b)} $x^2-(a+b)x+ab=0$
\item \textbf{c)} $x^2-2ax+a^2-b^2=0$
\item \textbf{d)} $ax^2-(a^2+1)x+a=0$, amb $a\neq0$
\end{itemize}
\end{minipage}\par\vspace{1.6mm}

\vspace{1mm}\hrule\vspace{3mm}

\section*{2. Equacions biquadrades i de grau superior}

\par\noindent\begin{minipage}{\linewidth}
\textbf{8.} Resol les equacions biquadrades següents:\par

\begin{tabular}{@{}p{.46\linewidth}p{.46\linewidth}@{}}
\textbf{a)} $x^4-13x^2+36=0$ & \textbf{b)} $x^4+3x^2-4=0$ \\[3pt]
\textbf{c)} $x^4+5x^2+6=0$ & \textbf{d)} $x^4-9x^2=0$ \\[3pt]
\end{tabular}\par
\end{minipage}\par\vspace{1.6mm}
\par\noindent\begin{minipage}{\linewidth}
\textbf{9.} Resol, reduint-les primer a una equació biquadrada:\par

\begin{itemize}[leftmargin=0pt,itemsep=6pt,topsep=4pt,parsep=0pt,label={}]
\item \textbf{a)} $(x^2-2)^2=1$
\item \textbf{b)} $(x^2+1)(x^2-3)=5$
\item \textbf{c)} $9x^4-6x^2+1=0$
\end{itemize}
\end{minipage}\par\vspace{1.6mm}
\par\noindent\begin{minipage}{\linewidth}
\textbf{10.} Resol factoritzant el polinomi amb Ruffini:\par

\begin{tabular}{@{}p{.46\linewidth}p{.46\linewidth}@{}}
\textbf{a)} $x^3-7x+6=0$ & \textbf{b)} $x^3-5x^2+7x-3=0$ \\[3pt]
\textbf{c)} $x^3-x^2+x-1=0$ &  \\[3pt]
\end{tabular}\par
\end{minipage}\par\vspace{1.6mm}
\par\noindent\begin{minipage}{\linewidth}
\textbf{11.} Resol les equacions següents:\par

\begin{itemize}[leftmargin=0pt,itemsep=6pt,topsep=4pt,parsep=0pt,label={}]
\item \textbf{a)} $x^4-2x^3-3x^2+8x-4=0$
\item \textbf{b)} $2x^4+x^3-8x^2-x+6=0$
\item \textbf{c)} $x^5-x^4-x+1=0$
\end{itemize}
\end{minipage}\par\vspace{1.6mm}
\par\noindent\begin{minipage}{\linewidth}
\textbf{12.} Troba el valor de $m$ perquè $x=2$ sigui solució de l'equació $x^3+mx^2-4x+4=0$ i, després, resol-la.\par
\end{minipage}\par\vspace{1.6mm}
\par\noindent\begin{minipage}{\linewidth}
\textbf{13.} Troba $k$ perquè $x=3$ sigui solució de l'equació $x^4+kx^2-18=0$ i, després, troba'n totes les solucions.\par
\end{minipage}\par\vspace{1.6mm}

\vspace{1mm}\hrule\vspace{3mm}

\section*{3. Equacions racionals}

\par\noindent\begin{minipage}{\linewidth}
\textbf{14.} Resol i comprova les solucions:\par

\begin{itemize}[leftmargin=0pt,itemsep=6pt,topsep=4pt,parsep=0pt,label={}]
\item \textbf{a)} $\dfrac{2}{x}+\dfrac{3}{x+1}=2$
\item \textbf{b)} $\dfrac{x}{x-2}=\dfrac{2}{x-2}+3$
\item \textbf{c)} $\dfrac{x}{x-3}+\dfrac{6}{x+3}=\dfrac{18}{x^2-9}$
\item \textbf{d)} $\dfrac{2x}{x+1}+\dfrac{1}{x}=3$
\end{itemize}
\end{minipage}\par\vspace{1.6mm}
\par\noindent\begin{minipage}{\linewidth}
\textbf{15.} Resol:\par

\begin{itemize}[leftmargin=0pt,itemsep=6pt,topsep=4pt,parsep=0pt,label={}]
\item \textbf{a)} $\dfrac{1}{x-1}-\dfrac{1}{x+1}=\dfrac{1}{12}$
\item \textbf{b)} $1+\dfrac{1}{x}=\dfrac{6}{x^2}$
\end{itemize}
\end{minipage}\par\vspace{1.6mm}
\par\noindent\begin{minipage}{\linewidth}
\textbf{16.} Troba el valor de $a$ perquè $x=3$ sigui solució de l'equació $\dfrac{a}{x-1}+\dfrac{1}{x+1}=1$ i, per a aquest valor, resol l'equació.\par
\end{minipage}\par\vspace{1.6mm}

\vspace{1mm}\hrule\vspace{3mm}

\section*{4. Equacions irracionals}

\par\noindent\begin{minipage}{\linewidth}
\textbf{17.} Resol i comprova les solucions:\par

\begin{tabular}{@{}p{.46\linewidth}p{.46\linewidth}@{}}
\textbf{a)} $\sqrt{x+5}=x-1$ & \textbf{b)} $\sqrt{2x+3}+x=6$ \\[3pt]
\textbf{c)} $x-\sqrt{x}=2$ & \textbf{d)} $2+\sqrt{x-1}=0$ \\[3pt]
\end{tabular}\par
\end{minipage}\par\vspace{1.6mm}
\par\noindent\begin{minipage}{\linewidth}
\textbf{18.} Resol les equacions amb dos radicals següents:\par

\begin{itemize}[leftmargin=0pt,itemsep=6pt,topsep=4pt,parsep=0pt,label={}]
\item \textbf{a)} $\sqrt{x+7}-\sqrt{x}=1$
\item \textbf{b)} $\sqrt{3x+1}-\sqrt{x-1}=2$
\item \textbf{c)} $\sqrt{x+4}+\sqrt{x-1}=5$
\item \textbf{d)} \emph{(més difícil)} $\sqrt{2x+1}+\sqrt{x-3}=2\sqrt{x}$
\end{itemize}
\end{minipage}\par\vspace{1.6mm}
\par\noindent\begin{minipage}{\linewidth}
\textbf{19.} Resol l'equació $\sqrt[3]{x^3-7}=x-1$. Cal comprovar les solucions quan l'arrel és cúbica?\par
\end{minipage}\par\vspace{1.6mm}
\par\noindent\begin{minipage}{\linewidth}
\textbf{20.} Troba $k$ perquè $x=5$ sigui solució de l'equació $\sqrt{x+k}=x-1$. Per a aquest valor de $k$, resol l'equació i comprova les solucions.\par
\end{minipage}\par\vspace{1.6mm}

\vspace{1mm}\hrule\vspace{3mm}

\section*{5. Equacions exponencials}

\par\noindent\begin{minipage}{\linewidth}
\textbf{21.} Resol expressant els dos membres com a potències de la mateixa base:\par

\begin{tabular}{@{}p{.46\linewidth}p{.46\linewidth}@{}}
\textbf{a)} $3^{x^2-1}=27$ & \textbf{b)} $5^{x^2-5x+6}=1$ \\[3pt]
\textbf{c)} $4^{x-1}=8^{x}$ & \textbf{d)} $\sqrt{2^{x+1}}=\dfrac{1}{4}$ \\[3pt]
\textbf{e)} $2^{x}\cdot 4^{x+1}=0{,}5$ &  \\[3pt]
\end{tabular}\par
\end{minipage}\par\vspace{1.6mm}
\par\noindent\begin{minipage}{\linewidth}
\textbf{22.} Resol traient factor comú:\par

\begin{itemize}[leftmargin=0pt,itemsep=6pt,topsep=4pt,parsep=0pt,label={}]
\item \textbf{a)} $2^{x+1}+2^{x}+2^{x-1}=28$
\item \textbf{b)} $3^{x+2}-3^{x}=72$
\end{itemize}
\end{minipage}\par\vspace{1.6mm}
\par\noindent\begin{minipage}{\linewidth}
\textbf{23.} Resol fent un canvi de variable:\par

\begin{itemize}[leftmargin=0pt,itemsep=6pt,topsep=4pt,parsep=0pt,label={}]
\item \textbf{a)} $4^{x}-5\cdot 2^{x}+4=0$
\item \textbf{b)} $9^{x}-2\cdot 3^{x}-3=0$
\item \textbf{c)} $2^{x}+2^{1-x}=3$
\item \textbf{d)} $a^{2x}-(a+1)\,a^{x}+a=0$, amb $a>0$ i $a\neq1$
\end{itemize}
\end{minipage}\par\vspace{1.6mm}
\par\noindent\begin{minipage}{\linewidth}
\textbf{24.} Resol aplicant logaritmes. Dona la solució exacta i una aproximació amb quatre decimals:\par

\begin{itemize}[leftmargin=0pt,itemsep=6pt,topsep=4pt,parsep=0pt,label={}]
\item \textbf{a)} $3^{x}=20$
\item \textbf{b)} $2^{x-1}=5$
\item \textbf{c)} $5e^{2x}=30$
\item \textbf{d)} \emph{(més difícil)} $2^{x}=3^{x-1}$
\end{itemize}
\end{minipage}\par\vspace{1.6mm}

\vspace{1mm}\hrule\vspace{3mm}

\section*{6. Equacions logarítmiques}

\par\noindent\begin{minipage}{\linewidth}
\textbf{25.} Resol i comprova les solucions:\par

\begin{tabular}{@{}p{.46\linewidth}p{.46\linewidth}@{}}
\textbf{a)} $\log(x^2-15)=1$ & \textbf{b)} $\log_3 x+\log_3(x-8)=2$ \\[3pt]
\textbf{c)} $2\log x=\log(5x-6)$ & \textbf{d)} $\log x+\log(x+3)=1$ \\[3pt]
\end{tabular}\par
\end{minipage}\par\vspace{1.6mm}
\par\noindent\begin{minipage}{\linewidth}
\textbf{26.} Resol aplicant les propietats dels logaritmes:\par

\begin{itemize}[leftmargin=0pt,itemsep=6pt,topsep=4pt,parsep=0pt,label={}]
\item \textbf{a)} $2\log x-\log(x-2)=2\log 3$
\item \textbf{b)} $\log(3x+1)-\log(2x-3)=1-\log 5$
\item \textbf{c)} $\ln x+\ln(x-1)=\ln 6$
\item \textbf{d)} $(\log x)^2-3\log x+2=0$
\item \textbf{e)} \emph{(més difícil)} $\log_2 x+\log_4 x=6$
\end{itemize}
\end{minipage}\par\vspace{1.6mm}
\par\noindent\begin{minipage}{\linewidth}
\textbf{27.} Troba $x$ aplicant la definició de logaritme:\par

\begin{tabular}{@{}p{.46\linewidth}p{.46\linewidth}@{}}
\textbf{a)} $\log_x 81=4$ & \textbf{b)} $\log_2(\log_3 x)=1$ \\[3pt]
\end{tabular}\par
\end{minipage}\par\vspace{1.6mm}
\par\noindent\begin{minipage}{\linewidth}
\textbf{28.} Aïlla $x$ en cada fórmula (totes les constants són positives i $b\neq1$):\par

\begin{itemize}[leftmargin=0pt,itemsep=6pt,topsep=4pt,parsep=0pt,label={}]
\item \textbf{a)} $y=a\cdot b^{x}$
\item \textbf{b)} $y=k\log\dfrac{x}{c}$
\item \textbf{c)} $\log x=\log a+2\log b-\dfrac12\log c$
\end{itemize}
\end{minipage}\par\vspace{1.6mm}

\vspace{1mm}\hrule\vspace{3mm}

\section*{7. Problemes}

\par\noindent\begin{minipage}{\linewidth}
\emph{\small En tots els problemes, comprova que les solucions obtingudes tenen sentit en el context.}\par
\end{minipage}\par\vspace{1.6mm}
\par\noindent\begin{minipage}{\linewidth}
\textbf{29.} Un rectangle té un perímetre de $34$ m i una diagonal de $13$ m. Quines dimensions té?\par
\end{minipage}\par\vspace{1.6mm}
\par\noindent\begin{minipage}{\linewidth}
\textbf{30.} Un grup d'amics lloga una furgoneta per $360$ € i se'n reparteixen el cost a parts iguals. Si fossin tres amics més, cadascun pagaria $20$ € menys. Quants amics són?\par
\end{minipage}\par\vspace{1.6mm}
\par\noindent\begin{minipage}{\linewidth}
\textbf{31.} Dues aixetes juntes omplen un dipòsit en $6$ hores. La primera, sola, tarda $5$ hores més que la segona. Quant tarda cada aixeta a omplir el dipòsit sola?\par
\end{minipage}\par\vspace{1.6mm}
\par\noindent\begin{minipage}{\linewidth}
\textbf{32.} Un ciclista recorre $120$ km. Si anés $10$ km/h més ràpid, tardaria una hora menys. A quina velocitat va?\par
\end{minipage}\par\vspace{1.6mm}
\par\noindent\begin{minipage}{\linewidth}
\textbf{33.} La diagonal d'un rectangle fa $5$ cm i la seva àrea és de $12\ \text{cm}^2$. Calcula'n les dimensions.\par
\end{minipage}\par\vspace{1.6mm}
\par\noindent\begin{minipage}{\linewidth}
\textbf{34.} \emph{(més difícil)} D'una làmina quadrada de cartró de $12$ cm de costat retallem un quadrat de costat $x$ a cada cantonada i, doblegant les vores, formem una caixa sense tapa. Quant ha de valer $x$ perquè el volum de la caixa sigui de $128\ \text{cm}^3$?\par
\end{minipage}\par\vspace{1.6mm}
\par\noindent\begin{minipage}{\linewidth}
\textbf{35.} El període d'un pèndol (el temps que tarda a fer una oscil·lació completa) és $T=2\pi\sqrt{\dfrac{L}{g}}$, amb $L$ en metres, $T$ en segons i $g=9{,}8\ \text{m/s}^2$.\par

\begin{itemize}[leftmargin=0pt,itemsep=6pt,topsep=4pt,parsep=0pt,label={}]
\item \textbf{a)} Quin és el període d'un pèndol d'$1$ m de longitud?
\item \textbf{b)} Quina longitud ha de tenir un pèndol perquè el seu període sigui de $3$ s?
\end{itemize}
\end{minipage}\par\vspace{1.6mm}
\par\noindent\begin{minipage}{\linewidth}
\textbf{36.} Una població de bacteris es duplica cada $3$ hores. Si inicialment n'hi ha $500$, el nombre de bacteris al cap de $t$ hores és $N(t)=500\cdot 2^{t/3}$. Quant de temps ha de passar perquè n'hi hagi $64\,000$?\par
\end{minipage}\par\vspace{1.6mm}
\par\noindent\begin{minipage}{\linewidth}
\textbf{37.} Un capital $C$ col·locat a interès compost del $r\,\%$ anual es converteix, al cap de $t$ anys, en $C_t=C\left(1+\dfrac{r}{100}\right)^{t}$.\par

\begin{itemize}[leftmargin=0pt,itemsep=6pt,topsep=4pt,parsep=0pt,label={}]
\item \textbf{a)} Demostra que el temps que tarda a duplicar-se no depèn de $C$ i troba'n una fórmula en funció de $r$.
\item \textbf{b)} Quant de temps tarda a duplicar-se al $3\,\%$? I al $6\,\%$?
\end{itemize}
\end{minipage}\par\vspace{1.6mm}
\par\noindent\begin{minipage}{\linewidth}
\textbf{38.} El cesi-137 té un període de semidesintegració de $30$ anys: cada $30$ anys, la quantitat que en queda es redueix a la meitat. Quant de temps ha de passar perquè en quedi només el $10\,\%$ de la quantitat inicial?\par
\end{minipage}\par\vspace{1.6mm}
\par\noindent\begin{minipage}{\linewidth}
\textbf{39.} El nivell sonor, en decibels, d'un so d'intensitat $I$ és $L=10\log\dfrac{I}{I_0}$, on $I_0$ és la intensitat mínima audible.\par

\begin{itemize}[leftmargin=0pt,itemsep=6pt,topsep=4pt,parsep=0pt,label={}]
\item \textbf{a)} Un concert fa $110$ dB i una conversa, $60$ dB. Quantes vegades més intens és el so del concert?
\item \textbf{b)} Si dupliquem la intensitat d'un so, quants decibels augmenta el nivell sonor?
\end{itemize}
\end{minipage}\par\vspace{1.6mm}
\par\noindent\begin{minipage}{\linewidth}
\textbf{40.} La magnitud $M$ d'un terratrèmol a l'escala de Richter és $M=\log\dfrac{A}{A_0}$, on $A$ és l'amplitud de les ones que registra el sismògraf i $A_0$ és una amplitud de referència. A més, l'energia $E$ alliberada (en joules) compleix $\log E=4{,}8+1{,}5M$.\par

\begin{itemize}[leftmargin=0pt,itemsep=6pt,topsep=4pt,parsep=0pt,label={}]
\item \textbf{a)} Quantes vegades més gran és l'amplitud d'un terratrèmol de magnitud $7$ que la d'un de magnitud $5$?
\item \textbf{b)} Quantes vegades més energia allibera el de magnitud $7$ que el de magnitud $5$?
\item \textbf{c)} Quina magnitud té un terratrèmol que allibera $1000$ vegades més energia que un de magnitud $4$?
\end{itemize}
\end{minipage}\par\vspace{1.6mm}
\par\noindent\begin{minipage}{\linewidth}
\textbf{41.} La pressió atmosfèrica disminueix amb l'altitud segons $p(h)=p_0\,e^{-kh}$, on $p_0=1013$ hPa és la pressió al nivell del mar, $h$ és l'altitud en metres i $k=1{,}2\cdot 10^{-4}\ \text{m}^{-1}$. Un altímetre baromètric mesura la pressió i en dedueix l'altitud.\par

\begin{itemize}[leftmargin=0pt,itemsep=6pt,topsep=4pt,parsep=0pt,label={}]
\item \textbf{a)} Quina pressió hi ha al cim de l'Aneto ($3404$ m)?
\item \textbf{b)} Aïlla $h$ en funció de $p$ (deixa-ho en funció de $p_0$ i $k$).
\item \textbf{c)} Si l'altímetre mesura una pressió de $700$ hPa, a quina altitud som?
\end{itemize}
\end{minipage}\par\vspace{1.6mm}
\vspace{1mm}\hrule\vspace{3mm}

\section*{Solucions}
\small\begin{multicols}{2}

\textbf{Bloc 1.}\par
\par\noindent\begin{minipage}{\linewidth}
\textbf{1.} a) $x=\dfrac{13}{2}$ · b) $x=\dfrac{15}{2}$ · c) qualsevol nombre real és solució (és una identitat) · d) no té solució.\par
\end{minipage}\par\vspace{1.6mm}
\par\noindent\begin{minipage}{\linewidth}
\textbf{2.} a) $x=\pm3$ · b) $x=0$, $x=-4$ · c) no té solució · d) $x=0$, $x=7$.\par
\end{minipage}\par\vspace{1.6mm}
\par\noindent\begin{minipage}{\linewidth}
\textbf{3.} a) $\Delta=1$; $x=2$, $x=3$ · b) $\Delta=0$; $x=\dfrac32$ (doble) · c) $\Delta=-3<0$; no té solució · d) $\Delta=49$; $x=\dfrac23$, $x=-\dfrac12$.\par
\end{minipage}\par\vspace{1.6mm}
\par\noindent\begin{minipage}{\linewidth}
\textbf{4.} a) $x=6$ (els termes de grau 2 s'anul·len) · b) $x=\pm1$ · c) $x=\dfrac{-1\pm\sqrt{41}}{2}$ · d) $x=5$, $x=-3$.\par
\end{minipage}\par\vspace{1.6mm}
\par\noindent\begin{minipage}{\linewidth}
\textbf{5.} a) $\Delta=36-4k=0 \Rightarrow k=9$; $x=3$ (doble) · b) $k=8$; l'altra solució és $x=4$ · c) $\Delta<0 \Leftrightarrow k>9$.\par
\end{minipage}\par\vspace{1.6mm}
\par\noindent\begin{minipage}{\linewidth}
\textbf{6.} a) $x^2+2x-15=0$ · b) $9x^2-12x+4=0$ (o qualsevol múltiple d'aquestes equacions).\par
\end{minipage}\par\vspace{1.6mm}
\par\noindent\begin{minipage}{\linewidth}
\textbf{7.} a) $x=1$ · b) $x=a$, $x=b$ · c) $x=a+b$, $x=a-b$ · d) $x=a$, $x=\dfrac1a$.\par
\end{minipage}\par\vspace{1.6mm}
\textbf{Bloc 2.}\par
\par\noindent\begin{minipage}{\linewidth}
\textbf{8.} a) $x=\pm2$, $x=\pm3$ · b) $x=\pm1$ ($t=-4$ no dona solucions) · c) no té solució ($t=-2$ i $t=-3$ són negatius) · d) $x=0$ (doble), $x=\pm3$.\par
\end{minipage}\par\vspace{1.6mm}
\par\noindent\begin{minipage}{\linewidth}
\textbf{9.} a) $x=\pm1$, $x=\pm\sqrt3$ · b) $x=\pm2$ · c) $x=\pm\dfrac{\sqrt3}{3}$ (totes dues dobles).\par
\end{minipage}\par\vspace{1.6mm}
\par\noindent\begin{minipage}{\linewidth}
\textbf{10.} a) $x=1$, $x=2$, $x=-3$ · b) $x=1$ (doble), $x=3$ · c) $x=1$ (el factor $x^2+1$ no té arrels reals).\par
\end{minipage}\par\vspace{1.6mm}
\par\noindent\begin{minipage}{\linewidth}
\textbf{11.} a) $x=1$ (doble), $x=\pm2$ · b) $x=\pm1$, $x=-2$, $x=\dfrac32$ · c) $x=1$ (doble), $x=-1$.\par
\end{minipage}\par\vspace{1.6mm}
\par\noindent\begin{minipage}{\linewidth}
\textbf{12.} $m=-1$; les solucions són $x=1$ i $x=\pm2$.\par
\end{minipage}\par\vspace{1.6mm}
\par\noindent\begin{minipage}{\linewidth}
\textbf{13.} $k=-7$; les solucions són $x=\pm3$ ($t=-2$ no dona solucions).\par
\end{minipage}\par\vspace{1.6mm}
\textbf{Bloc 3.}\par
\par\noindent\begin{minipage}{\linewidth}
\textbf{14.} a) $x=2$, $x=-\dfrac12$ · b) no té solució ($x=2$ anul·la el denominador) · c) $x=-12$ ($x=3$ no és vàlida) · d) $x=-1\pm\sqrt2$.\par
\end{minipage}\par\vspace{1.6mm}
\par\noindent\begin{minipage}{\linewidth}
\textbf{15.} a) $x=\pm5$ · b) $x=2$, $x=-3$.\par
\end{minipage}\par\vspace{1.6mm}
\par\noindent\begin{minipage}{\linewidth}
\textbf{16.} $a=\dfrac32$; les solucions són $x=3$ i $x=-\dfrac12$.\par
\end{minipage}\par\vspace{1.6mm}
\textbf{Bloc 4.}\par
\par\noindent\begin{minipage}{\linewidth}
\textbf{17.} a) $x=4$ ($x=-1$ no és vàlida) · b) $x=3$ ($x=11$ no és vàlida) · c) $x=4$ ($x=1$ no és vàlida) · d) no té solució ($x=5$ no és vàlida: una arrel quadrada no pot valer $-2$).\par
\end{minipage}\par\vspace{1.6mm}
\par\noindent\begin{minipage}{\linewidth}
\textbf{18.} a) $x=9$ · b) $x=1$, $x=5$ · c) $x=5$ · d) $x=4$ ($x=-\dfrac47$ no és vàlida).\par
\end{minipage}\par\vspace{1.6mm}
\par\noindent\begin{minipage}{\linewidth}
\textbf{19.} $x=2$, $x=-1$. En elevar al cub no apareixen solucions falses, però sempre és bo comprovar-les.\par
\end{minipage}\par\vspace{1.6mm}
\par\noindent\begin{minipage}{\linewidth}
\textbf{20.} $k=11$; l'única solució és $x=5$ ($x=-2$ no és vàlida).\par
\end{minipage}\par\vspace{1.6mm}
\textbf{Bloc 5.}\par
\par\noindent\begin{minipage}{\linewidth}
\textbf{21.} a) $x=\pm2$ · b) $x=2$, $x=3$ · c) $x=-2$ · d) $x=-5$ · e) $x=-1$.\par
\end{minipage}\par\vspace{1.6mm}
\par\noindent\begin{minipage}{\linewidth}
\textbf{22.} a) $x=3$ · b) $x=2$.\par
\end{minipage}\par\vspace{1.6mm}
\par\noindent\begin{minipage}{\linewidth}
\textbf{23.} a) $x=0$, $x=2$ · b) $x=1$ ($3^x=-1$ no té solució) · c) $x=0$, $x=1$ · d) $x=0$, $x=1$ (per a qualsevol valor de $a$).\par
\end{minipage}\par\vspace{1.6mm}
\par\noindent\begin{minipage}{\linewidth}
\textbf{24.} a) $x=\log_3 20=\dfrac{\log 20}{\log 3}\approx 2{,}7268$ · b) $x=1+\log_2 5\approx 3{,}3219$ · c) $x=\dfrac{\ln 6}{2}\approx 0{,}8959$ · d) $x=\dfrac{\log 3}{\log 3-\log 2}\approx 2{,}7095$.\par
\end{minipage}\par\vspace{1.6mm}
\textbf{Bloc 6.}\par
\par\noindent\begin{minipage}{\linewidth}
\textbf{25.} a) $x=\pm5$ (totes dues vàlides) · b) $x=9$ ($x=-1$ no és vàlida) · c) $x=2$, $x=3$ · d) $x=2$ ($x=-5$ no és vàlida).\par
\end{minipage}\par\vspace{1.6mm}
\par\noindent\begin{minipage}{\linewidth}
\textbf{26.} a) $x=3$, $x=6$ · b) $x=7$ · c) $x=3$ ($x=-2$ no és vàlida) · d) $x=10$, $x=100$ · e) $x=16$.\par
\end{minipage}\par\vspace{1.6mm}
\par\noindent\begin{minipage}{\linewidth}
\textbf{27.} a) $x=3$ ($x=-3$ no pot ser base) · b) $x=9$.\par
\end{minipage}\par\vspace{1.6mm}
\par\noindent\begin{minipage}{\linewidth}
\textbf{28.} a) $x=\log_b\dfrac{y}{a}=\dfrac{\log y-\log a}{\log b}$ · b) $x=c\cdot 10^{y/k}$ · c) $x=\dfrac{ab^2}{\sqrt{c}}$.\par
\end{minipage}\par\vspace{1.6mm}
\textbf{Bloc 7.}\par
\par\noindent\begin{minipage}{\linewidth}
\textbf{29.} $12$ m i $5$ m.\par
\end{minipage}\par\vspace{1.6mm}
\par\noindent\begin{minipage}{\linewidth}
\textbf{30.} $6$ amics (paguen $60$ € cadascun).\par
\end{minipage}\par\vspace{1.6mm}
\par\noindent\begin{minipage}{\linewidth}
\textbf{31.} La segona, $10$ h; la primera, $15$ h.\par
\end{minipage}\par\vspace{1.6mm}
\par\noindent\begin{minipage}{\linewidth}
\textbf{32.} $30$ km/h.\par
\end{minipage}\par\vspace{1.6mm}
\par\noindent\begin{minipage}{\linewidth}
\textbf{33.} $4$ cm i $3$ cm (surt l'equació biquadrada $x^4-25x^2+144=0$).\par
\end{minipage}\par\vspace{1.6mm}
\par\noindent\begin{minipage}{\linewidth}
\textbf{34.} $x=2$ cm. L'equació $x(12-2x)^2=128$ també té la solució $x=8$, que no té sentit ($12-2x<0$).\par
\end{minipage}\par\vspace{1.6mm}
\par\noindent\begin{minipage}{\linewidth}
\textbf{35.} a) $T\approx 2{,}01$ s · b) $L=\dfrac{9{,}8\cdot 9}{4\pi^2}\approx 2{,}23$ m.\par
\end{minipage}\par\vspace{1.6mm}
\par\noindent\begin{minipage}{\linewidth}
\textbf{36.} $21$ hores.\par
\end{minipage}\par\vspace{1.6mm}
\par\noindent\begin{minipage}{\linewidth}
\textbf{37.} a) $t=\dfrac{\log 2}{\log\left(1+\frac{r}{100}\right)}$ · b) $t\approx 23{,}45$ anys al $3\,\%$ i $t\approx 11{,}90$ anys al $6\,\%$.\par
\end{minipage}\par\vspace{1.6mm}
\par\noindent\begin{minipage}{\linewidth}
\textbf{38.} $\left(\dfrac12\right)^{t/30}=0{,}1 \Rightarrow t=\dfrac{30}{\log 2}\approx 99{,}66$ anys.\par
\end{minipage}\par\vspace{1.6mm}
\par\noindent\begin{minipage}{\linewidth}
\textbf{39.} a) $10^5=100\,000$ vegades · b) $10\log 2\approx 3{,}01$ dB.\par
\end{minipage}\par\vspace{1.6mm}
\par\noindent\begin{minipage}{\linewidth}
\textbf{40.} a) $10^2=100$ vegades · b) $10^{3}=1000$ vegades · c) $M=6$.\par
\end{minipage}\par\vspace{1.6mm}
\par\noindent\begin{minipage}{\linewidth}
\textbf{41.} a) $p\approx 673{,}3$ hPa · b) $h=\dfrac{1}{k}\ln\dfrac{p_0}{p}$ · c) $h\approx 3080$ m.\par
\end{minipage}\par\vspace{1.6mm}
\end{multicols}
\end{document}
